\begin{textsample}{NEG Dim 4 – 2014\_c5 – Score 1.00 – 2014\_c5/enade\_2014\_hist\_b11f5fba33.txt}  \label{ex:f4_neg_200}
CAPÍTULO 5 DISTRIBUIÇÃO DOS CONCEITOS 5.1 PANORAMA NACIONAL DA DISTRIBUIÇÃO DOS CONCEITOS O Conceito Enade é um indicador de qualidade que avalia o desempenho dos estudantes a partir dos resultados do ENADE.
Os valores possíveis vão de 1 ( pior situação ) a 5 ( melhor situação ).
As Unidades de Observação com apenas um ou sem nenhum concluinte participante não obtêm o Conceito Enade, ficando Sem Conceito ( SC ).
A Tabela 5.1 apresenta a quantidade e distribuição de cursos de História ( Bacharelado ) participantes do ENADE/2014, por faixa de conceito e Grande Região.
Mais uma vez, ressaltando que a diferença entre os cursos tabulados neste capítulo e no capítulo 2 são os cursos sem conceito, em princípio, aqueles sem alunos concluintes que participassem da prova ou que tivessem tão somente um aluno concluinte²².
Observando -se os dados da Tabela 5.1, nota -se que, dos 51 cursos participantes, 19 ( 37,3 \% ) classificaram -se com conceito 3, o valor modal.
Este foi também o conceito modal em mais três regiões : Nordeste ( 44,4 \% ), Sudeste ( 36,4 \% ) e Sul ( 45,5 \% ).
Na região Centro-Oeste, a moda foi o conceito 2 ( 50,0 \% ).
O conceito 2 foi o segundo mais frequente em nível nacional ( 25,5 \%, correspondendo a 13 cursos ) e os conceito 1 e 4, os terceiros ( 11,8 \%, cada ), correspondendo a seis cursos cada.
Houve, ainda, três cursos ( 5,9 \% ) que receberam conceito 5.
Dos 51 cursos de História ( Bacharelado ), apenas quatro cursos ( 7,8 \% ) ficaram sem conceito ( SC ). ²² Estes cursos com somente um concluinte não foram considerados no capítulo 2.
A região Norte participou com cinco cursos ou 9,8 \% do total nacional.
Nesta região, o conceito 1 foi atribuído a três cursos dos cinco participantes, o que equivale a 60,0 \% do total regional, além de ser o conceito modal na região.
Nenhum curso ficou sem conceito ( SC ) ou com conceitos 4 ou 5.
Dois cursos foram avaliados com conceitos 2 e 3 ( dois cursos, cada ), correspondendo a 20,0 \%, cada.
A região Nordeste participou com nove cursos ou 17,6 \% do total nacional.
Nesta região, quatro cursos, 44,4 \% em termos regionais, obtiveram conceito 3, o conceito modal para a região.
Os conceitos 2 e 4 foram atribuídos a, respectivamente, dois e um curso ( 22,2 \% e 11,1 \% ).
Dois cursos ( 22,2 \% ) obtiveram conceito 1.
Nenhum curso ficou sem conceito ou com conceito 5 nessa região.
Dos 22 cursos participantes da região Sudeste ( 43,1 \% do total nacional ), o conceito 3 também foi o conceito modal, e foi obtido por oito ( 36,4 \% ) cursos.
O conceito 1 foi atribuído a um curso ( 4,5 \% ) e os conceitos 2 e 4, com quatro e cinco cursos, correspondendo a 18,2 \% e 22,7 \%, respectivamente.
O conceito 5 foi atribuído a três cursos ( 13,6 \% ).
E um curso ( 4,5 \% ) ficou sem conceito ( SC ).
Dos 11 cursos da região Sul, cinco cursos, 45,5 \% em termos regionais, obtiveram conceito 3 e quatro cursos receberam conceito 2.
Dois cursos ficaram sem conceito ( SC ) e nenhum curso foi atribuído aos demais conceitos nesta região.
Já dos quatro cursos participantes na região Centro-Oeste, dois receberam conceito 2, o conceito modal ( 50,0 \% ).
Nesta região, um curso ficou sem conceito ( SC ) e um curso foi avaliado com conceito 3.
Nenhum curso foi avaliado nos demais conceitos nesta região.
O Gráfico 5.1 apresenta a distribuição do Conceito Enade segundo Grande Região : poligonais mais à esquerda representam Grandes Regiões com pior distribuição, e poligonais mais à direita, Grandes Regiões com melhores conceitos.
A região Norte ( linha azul ) apresenta os cursos em IES com conceitos mais baixos, ao passo que a região Sudeste ( linha vermelha ) apresentou a distribuição com valores maiores.
Cursos em IES das regiões Centro-Oeste ( linha preta ), Sul ( linha roxa ) e Nordeste ( linha verde ) apresentam uma situação intermediária.
A Tabela 5.2 apresenta a quantidade e distribuição de cursos de História ( Licenciatura ) participantes do ENADE/2014, por faixa de conceito e Grande Região.
A diferença entre os cursos tabulados neste capítulo e no capítulo 2 são os cursos sem conceito ( SC ), em princípio, aqueles sem alunos concluintes que participassem da prova ou que tivessem tão somente um aluno concluinte²³.
Observando -se os dados da Tabela 5.2, nota -se que, dos 364 cursos participantes, 143 ( 39,3 \% ) classificaram -se com conceito 3, o valor modal, assim como em História ( Bacharelado ).
Este foi também o conceito modal em mais duas regiões : Sudeste ( 52,5 \% ) e Sul ( 43,1 \% ).
O conceito 2 foi o segundo mais frequente em nível nacional ( 32,7 \%, correspondendo a 119 cursos ) e o conceito 4, o terceiro ( 13,5 \%, correspondendo a 49 cursos ).
Houve, ainda, 14 cursos ( 3,8 \% ) que receberam conceito 5 e outros 27 cursos que receberam conceito 1.
Dos 364 cursos de História ( Licenciatura ), 12 cursos ficaram sem conceito ( SC ). ²³ Estes cursos com somente um concluinte não foram considerados no capítulo 2.
A região Norte participou com 22 cursos ou 6,0 \% do total nacional.
Nesta região, o conceito 2 foi atribuído a 11 cursos dos 22 participantes, o que equivale a 50,0 \% do total regional, além de ser o conceito modal na região.
Um curso ficou sem conceito ( SC ).
Dois cursos foram avaliados com conceito 1 e oito com conceito 3.
Nenhum curso foi avaliado com conceito 4 ou 5.
A região Nordeste participou com 115 cursos ou 31,6 \% do total nacional.
Nesta região, 54 cursos, 47,0 \% em termos regionais, obtiveram conceito 2, também o conceito modal para a região.
Os conceitos 1 e 3 foram atribuídos a, respectivamente, 18 e 31 cursos ( 15,7 \% e 27,0 \% ).
Enquanto oito cursos obtiveram conceito 4 e um curso obteve conceito 5.
Três cursos ficaram sem conceito ( SC ).
Dos 120 cursos participantes da região Sudeste ( 33,0 \% do total nacional ), o conceito 3, conceito modal, foi obtido por 63 ( 52,5 \% ) cursos.
O conceito 1 foi atribuído a três cursos ( 2,5 \% ) e os conceitos 2 e 4, a 25 ( 20,8 \% ) e 18 ( 15,0 \% ) cursos, respectivamente.
O conceito 5 foi atribuído a nove cursos ( 7,5 \% ).
Somente dois cursos ( 1,7 \% ) ficou sem conceito ( SC ).
Os 65 cursos da região Sul ficaram distribuídos em todas as faixas de conceitos.
A predominância, mais uma vez, foi do conceito 3 com 43,1 \% ( conceito modal ), ou seja, 28 dos 65 cursos participantes na região Sul.
O conceito 2 foi atribuído a nove cursos ( 13,8 \% ), o conceito 1 a dois ( 3,1 \% ) e o conceito 4, a 18 cursos ( 27,7 \% ).
O conceito 5 foi recebido por quatro cursos ( 6,2 \% ), enquanto quatro cursos ( 6,2 \% ), na região Sul, ficaram sem conceito ( SC ).
Já dos 42 cursos participantes na região Centro-Oeste, 20 receberam conceito 2, o conceito modal ( 47,6 \% ).
Nesta região, dois cursos ficaram sem conceito ( SC ) e nenhum curso obteve conceito 5.
Já, dois ( 4,8 \% ) cursos foram \textbf{alocados} no conceito 1, enquanto 13 ( 31,0 \% ) dos cursos obtiveram conceito 3 e cinco ( 11,9 \% ) foram avaliados com conceito 4.
O Gráfico 5.2 apresenta a distribuição do Conceito Enade segundo Grande Região : poligonais mais à esquerda representam Grandes Regiões com pior distribuição, e poligonais mais à direita, Grandes Regiões com melhores conceitos.
As regiões Norte ( linha azul ) e Nordeste ( linha verde ) apresentam os cursos em IES com conceitos mais baixos, ao passo que as regiões Sudeste ( linha vermelha ) e Sul ( linha roxa ) apresentam as distribuições com valores maiores. 5.2 CONCEITOS POR CATEGORIA ADMINISTRATIVA E POR GRANDE REGIÃO A Tabela 5.3 apresenta a distribuição dos cursos participantes do ENADE/2014 de História ( Bacharelado ), por Categoria Administrativa, de acordo com os conceitos por eles alcançados, segundo as Grandes Regiões brasileiras.
Dos 51 cursos participantes, 40 ( 78,4 \% ) eram ministrados em Instituições Públicas e 11 ( 21,6 \% ), em Privadas.
De acordo com as informações da Tabela 5.3, em termos nacionais, entre Instituições Públicas e Privadas : das três IES que receberam o conceito 5, duas eram Públicas e uma Privada.
Dos 40 cursos participantes de IES Públicas, o conceito 3 foi o valor modal, atribuído a 14 cursos ( 35,0 \% ).
Entre os demais cursos de Instituições Públicas participantes, cinco cursos ( 12,5 \% ) foram avaliados com conceito 1, 12 cursos obtiveram conceito 2 ( 30,0 \% da categoria ), cinco cursos receberam conceito 4 ( 12,5 \% ) e, como já comentado, dois foram \textbf{alocados} ao conceito 5 ( 5,0 \% ).
Nesta categoria, somente dois cursos ficaram sem conceito ( SC ).
Na rede privada, o conceito modal foi 3, com cinco cursos dos 11 da categoria.
Entre os demais cursos participantes, quatro ficaram com os demais conceitos restantes ( um curso, cada ).
Dois cursos ficaram sem conceito ( SC ).
Na análise por região, observa -se que, na região Norte, as Instituições Privadas não tiveram participação.
As Instituições Públicas participaram com os cinco cursos na região Norte ( 100,0 \% do total regional ), dos quais três obtiveram conceitos 1, o valor modal, seguido por dois cursos que receberam os conceitos 2 e 3 ( um curso, cada ).
Nenhum curso ficou sem conceito ( SC ) ou recebeu conceitos 4 ou 5.
Na região Nordeste, a rede privada concentrou apenas um dos nove cursos participantes, o equivalente a 11,1 \% do total da região.
Este curso foi avaliado com conceito 1.
As Instituições Públicas dessa região participaram com oito cursos ( 88,9 \% ), dos quais quatro ( 50,0 \% ) obtiveram conceito 3, o valor modal, seguido por dois cursos que receberam conceito 2 ( 25,0 \% ).
Dois cursos foram \textbf{alocados} aos conceitos 1 e 4, enquanto nenhum curso obteve conceito 5 ou ficou sem conceito ( SC ).
Na região Sudeste, a proporção de cursos da rede privada correspondeu a 27,3 \%, com seis dos 22 cursos participantes.
Nesta categoria, na região Sudeste, o conceito modal foi 3 ( três cursos ).
Nesta combinação de Categoria Administrativa e Grande Região, um curso não recebeu conceito ( SC ) e outros dois receberam conceitos 4 e 5 ( um curso, cada ).
Nenhum curso obteve conceitos 1 e 2.
Entre os 16 cursos oferecidos por Instituições Públicas na região Sudeste, o conceito modal foi 3 ( cinco cursos ), e os demais cursos receberam os conceitos 1 ( um curso ), 2 ( quatro cursos ), 4 ( quatro cursos ) e 5 ( dois cursos ).
Na região Sudeste, nesta categoria, nenhum curso ficou sem conceito ( SC ).
As Instituições Privadas concentraram quatro dos 11 cursos participantes da região Sul, 36,4 \% do total regional, maior proporção de Instituições Privadas se comparada com as demais regiões.
Desses, dois ficaram com conceito 3 e um com conceito 2.
Somente um curso ficou sem conceito ( SC ).
Nenhum curso foi avaliado nos demais conceitos restante nesta categoria e região.
As Instituições Públicas na região Sul participaram com sete cursos ( 63,6 \% ), a menor proporção desta categoria dentre todas as regiões, sendo seis cursos avaliados com conceitos 2 e 3 ( três cursos cada, conceitos modais ) e um curso que ficou sem conceito.
Nenhum curso obteve os demais conceitos restantes nesta categoria e região.
Na região Centro-Oeste, as Instituições Privadas não tiveram participação.
Dos quatro cursos oferecidos por Instituições Públicas, dois foram avaliados no conceito 2, conceito modal, seguidos de um curso que obteve conceito 3.
Um curso ficou sem conceito ( SC ) e nenhum curso foi avaliado nos demais conceitos nesta categoria e região.
O Gráfico 5.3 apresenta a distribuição do Conceito Enade segundo a Categoria Administrativa da IES.
Os cursos em IES Privadas ( linha verde ) apresentam uma poligonal mais à direita do que os em IES Públicas ( linha azul ) e, consequentemente, uma distribuição de Conceitos Enade com valores maiores.
A Tabela 5.4 apresenta a distribuição dos cursos participantes do ENADE/2014 de História ( Licenciatura ), por Categoria Administrativa, de acordo com os conceitos por eles alcançados, segundo as Grandes Regiões brasileiras.
Dos 364 cursos participantes, 198 ( 54,4 \% ) eram ministrados em Instituições Públicas e 166 ( 45,6 \% ), em Privadas.
De acordo com as informações da Tabela 5.4, em termos nacionais, entre Instituições Públicas e Privadas : das 14 IES que receberam o conceito 5, 11 eram Públicas e três Privadas.
Dos 198 cursos participantes de IES Públicas, o conceito 2 foi o valor modal, atribuído a 71 cursos ( 35,9 \% ).
Entre os demais cursos de Instituições Públicas participantes, somente quatro cursos não obtiveram conceito ( SC ), 62 cursos receberam conceito 4 ( 31,3 \% ) e, como já comentado, 11 foram \textbf{alocados} ao conceito 5 ( 5,6 \% ).
Nesta categoria, 23 ( 11,6 \% ) cursos obtiveram conceito 1, enquanto 27 cursos, obtiveram conceito 4.
Na rede privada, o conceito modal também foi 3, com 81 cursos dos 166 da categoria.
Entre os cursos participantes, quatro receberam conceito 1 e 48, conceito 2.
O conceito 4 foi atribuído a 22 cursos e o conceito 5, a três cursos.
Nesta Categoria Administrativa, oito cursos ficaram sem conceito ( SC ).
Na análise por região, observa -se que, na região Norte, as Instituições Privadas participaram com sete cursos ( 31,8 \% do total regional ), dos quais um ficou sem conceito ( SC ) e nenhum obteve os conceitos 1, 4 ou 5.
Os conceitos modais para as Instituições Privadas na região foram 2 e 3 ( três cursos, cada ).
As Instituições Públicas participaram com 15 cursos na região Norte ( 68,2 \% do total regional ).
Desses, oito cursos obtiveram conceito 2, conceito modal nessa região.
Cinco cursos receberam conceito 3, enquanto dois cursos obtiveram conceito 1.
Nenhum curso ficou sem conceito ( SC ) ou com conceitos 4 ou 5 nessa categoria.
Na região Nordeste, a rede privada concentrou 18 dos 115 cursos participantes, o equivalente a 15,7 \% do total da região, a menor proporção desta categoria dentre todas as regiões.
Dentre os cursos oferecidos por IES Privadas no Nordeste, predominaram os cursos com conceito 2, 13 cursos correspondendo a 72,2 ( conceito modal ).
Já dois cursos receberam o conceito 1 ( 11,1 \% ), enquanto o conceito 3 foi atribuído a três cursos ( 16,7 \% ).
Nenhum curso ficou sem conceito ( SC ) ou obteve conceitos 4 ou 5.
As Instituições Públicas dessa região participaram com 97 cursos ( 84,3 \% ), a maior proporção desta categoria dentre todas as regiões, dos quais 41 cursos ( 42,3 \% ) obtiveram conceitos 2, também o valor modal nessa combinação de Categoria Administrativa e Grande Região, seguido por 28 cursos ( 28,9 \% ) que obtiveram conceito 3.
Já 16 cursos ( 16,5 \% ) foram \textbf{alocados} ao conceito 1 e oito ( 8,2 \% ) no conceito 4.
Somente um curso ( 1 \% ) recebeu conceito 5 e três cursos ( 3,1 \% ) ficaram sem conceito ( SC ) nesta combinação de categoria e região.
Na região Sudeste, a proporção de cursos da rede privada, 75,0 \%, foi mais elevada do que nas demais regiões brasileiras, correspondendo a 90 dos 120 cursos participantes.
Nesta categoria, na região Sudeste, o conceito modal foi 3 ( 51 cursos ).
Nesta combinação de Categoria Administrativa e Grande Região, somente um curso não recebeu conceito ( SC ).
Os demais foram avaliados com conceito 1 ( um curso ), conceito 2 ( 21 cursos ), conceito 4 ( 15 cursos ) e conceito 5 ( um curso ).
Entre os 30 cursos oferecidos por Instituições Públicas na região Sudeste, o conceito modal foi 3 ( 12 cursos ), e os demais cursos receberam os conceitos 2 ( quatro cursos ) e 1 ( dois cursos ).
Na região Sudeste, nesta categoria, três cursos foram avaliados com conceitos 4 e oito com conceito 5.
Um curso ficou sem conceito ( SC ).
As Instituições Privadas concentraram 39 dos 65 cursos participantes da região Sul, 60,0 \% do total regional.
Desses, 19 ficaram com conceito 3, o conceito modal.
Nesta combinação de Categoria Administrativa e Grande Região, sete cursos ficaram com conceito 4, seis cursos obtiveram conceito 2.
Dois cursos receberam conceito 5, enquanto quatro cursos ficaram sem conceito ( SC ).
Um curso obteve conceito 1.
As Instituições Públicas na região Sul participaram com 26 cursos ( 40,0 \% ), e o conceito modal foi 4, com 11 cursos.
Os demais cursos foram avaliados com conceitos 1 ( um curso ), 2 ( três cursos ), 3 ( nove cursos ) e 5 ( dois cursos ).
Nenhum curso ficou sem conceito ( SC ).
Na região Centro-Oeste, 12 dos 42 cursos participantes eram de Instituições Privadas ( 28,6 \% em termos regionais ).
Destes, dez concentraram -se nos conceitos 2 e 3 ( cinco cursos, cada ), conceitos modais.
Dois cursos ficaram sem conceito ( SC ) e nenhum curso foi avaliado nos demais conceitos nesta categoria e região.
Dos 30 cursos oferecidos por Instituições Públicas, 15 foram avaliados no conceito 2, conceito modal, e dois cursos foram \textbf{alocados} no conceito 1, e oito, no conceito 3.
Cinco cursos foram avaliados com conceito 4 e nenhum curso de IES pública ficou sem conceito ( SC ) ou recebeu conceito 5 nessa categoria e região.
O Gráfico 5.4 apresenta a distribuição do Conceito Enade segundo a Categoria Administrativa da IES.
Como as linhas se cruzam, fica patente que as IES Privadas ( linha verde ) apresentam uma distribuição de Conceitos Enade mais concentrada do que as IES Públicas ( linha azul ). 5.3 CONCEITOS POR ORGANIZAÇÃO ACADÊMICA E POR GRANDE REGIÃO Na Tabela 5.5, encontra -se a distribuição dos conceitos atribuídos aos cursos participantes do ENADE/2014 na Área de História ( Bacharelado ), por Organização Acadêmica, segundo as Grandes Regiões brasileiras.
Dos 51 cursos de História ( Bacharelado ) participantes, 50 eram oferecidos em Universidades, apenas um em Centros Universitários e nenhum em Faculdades.
Esta distribuição corresponde a, respectivamente, 98,0 \%, 2,0 \% e 0,0 \% dos cursos.
De acordo com os dados apresentados, dos três cursos avaliados com conceito 5, todos eram vinculados a Universidades.
Somente quatro cursos oferecidos em Universidades ficaram sem conceito ( SC ).
Esse tipo de Organização Acadêmica teve o conceito 3 como modal, com 18 cursos ( 36,0 \% ).
Os demais cursos de Universidades avaliados receberam os conceitos 1 ( seis cursos ), 2 ( 13 cursos ), 4 ( seis cursos ) e conceito 5 ( três cursos, como já mencionado ).
Os Centros Universitários participaram com apenas um curso, este curso obteve conceito 3.
As Faculdades não tiveram participação, como já dito anteriormente.
Considerando -se separadamente as regiões brasileiras, verifica -se que, na região Norte, as Universidades concentraram todos os cinco cursos participantes nesta região.
Três cursos ficaram \textbf{alocados} ao conceito 1, conceito modal, seguidos de dois cursos que obtiveram conceitos 2 e 3 ( um curso, cada ).
Nenhum curso ficou sem conceito ( SC ) ou recebeu os demais conceitos nesta organização e região.
Os Centros Universitários e as Faculdades não tiveram participação na região Norte.
Na região Nordeste, as Universidades participaram com todos os nove cursos na Área de História ( Bacharelado ), nesta região.
O conceito modal foi 3, com quatro cursos, enquanto o conceito 4 foi atribuído a um curso.
Quatro cursos obtiveram conceitos 1 e 2 ( dois cursos, cada ).
Nenhum dos cursos oferecidos por Universidades no Nordeste ficou sem conceito ( SC ) ou ficou com conceito 5.
Os Centros Universitários e as Faculdades não tiveram participação na região Nordeste.
Na região Sudeste, as Universidades concentraram todos os 22 cursos participantes nesta região.
Entre os cursos em Universidades, nessa região, o conceito modal foi 3 com oito cursos, enquanto um curso ficou sem conceito ( SC ).
Os demais cursos receberam os conceitos 1 ( um curso ), 2 ( quatro cursos ), 4 ( cinco cursos ) e 5 ( três cursos ).
Os Centros Universitários e as Faculdades não tiveram participação na região Sudeste.
Dos 11 cursos da região Sul, dez eram de Universidades, para os quais os conceitos modais foi 2 e 3, com quatro cursos cada.
Nesse tipo de organização, dois cursos ficaram sem conceito ( SC ).
Os demais cursos não receberam os demais conceitos restantes nesta organização e região.
Os Centros Universitários participaram com somente um curso na região Sul, este curso ficou conceito 3.
As Faculdades não tiveram participação na região Sul.
Na região Centro-Oeste, dos quatro cursos participantes, todos eram de Universidades.
Nesse tipo de organização, o conceito modal foi 2, com dois cursos seguidos de um curso que obteve conceito 3.
Um curso ficou sem conceito ( SC ) e nenhum curso foi avaliados com os demais conceitos restantes nesta organização e região.
Os Centros Universitários e as Faculdades não tiveram participação na região Centro-Oeste.
O Gráfico 5.5 apresenta a distribuição do Conceito Enade segundo a Organização Acadêmica da IES.
Como as linhas se cruzam, fica patente que os cursos em Centros universitários ( linha verde ) apresentam uma distribuição de Conceitos Enade mais concentrada do que os cursos em Universidades ( linha azul ).
Na Tabela 5.6, encontra -se a distribuição dos conceitos atribuídos aos cursos participantes do ENADE/2014 na Área de História ( Licenciatura ), por Organização Acadêmica, segundo as Grandes Regiões brasileiras.
Dos 364 cursos de História ( Licenciatura ) participantes, 246 eram oferecidos em Universidades, 82 em Faculdades e os demais 36, em Centros Universitários.
Esta distribuição corresponde a, respectivamente, 67,6 \%, 22,5 \% e 9,9 \% dos cursos.
De acordo com os dados apresentados, dos 14 cursos avaliados com conceito 5, 13 eram vinculados a Universidades.
Apenas oito cursos oferecidos em Universidades ficaram sem conceito ( SC ).
Esse tipo de Organização Acadêmica teve o conceito 3 como modal, com 94 cursos ( 38,2 \% ).
Os demais cursos de Universidades avaliados receberam os conceitos 1 ( 16 cursos ), 2 ( 79 cursos ), 4 ( 36 cursos ) e conceito 5 ( 13 cursos, como já mencionado ).
Entre os cursos em Centros Universitários, o conceito modal foi 3, com 18 cursos ( 50,0 \% ).
Neste tipo de Organização Acadêmica, os cursos receberam os conceitos 1 ( um curso ), 2 ( seis cursos ) e conceito 4 ( nove cursos ).
Nenhum curso ficou com conceito 5 e nenhum curso ficou sem conceito ( SC ).
Nas Faculdades, dois dos 82 cursos ficaram sem conceito ( SC ), 34 ( 41,5 \% ) ficaram \textbf{alocados} no conceito 2, conceito modal, e 31 cursos receberam conceito 3.
Dos demais cursos neste tipo de Organização Acadêmica, dez receberam conceito 1, quatro, conceito 4, e um recebeu conceito 5.
Considerando -se separadamente as regiões brasileiras, verifica -se que, na região Norte, as Universidades concentraram 16 dos 22 cursos participantes, dos quais nove cursos obtiveram conceito 2, o conceito modal, seguido de cinco cursos com conceito 3 e dois cursos que receberam conceito 1.
Nenhum curso ficou \textbf{alocado} aos conceitos 4 ou 5 e nenhum curso ficou sem conceito ( SC ) nesta categoria e região.
Os Centros Universitários da região Norte foram representados por somente um curso, este curso foi avaliado com conceito 2.
As Faculdades participaram com cinco cursos na região Norte.
Um curso ficou sem conceito ( SC ) e outro ficou com conceito 2.
Os cursos oferecidos por Faculdades receberam conceito 3 ( três cursos ), conceito modal.
Nenhum curso foi \textbf{alocado} nos demais conceitos nesta organização e região.
Na região Nordeste, as Universidades participaram com 92 dos 115 cursos na Área de História ( Licenciatura ).
O conceito modal foi 2, com 41 cursos, enquanto os demais cursos oferecidos por Universidades receberam conceitos 1 ( dez cursos ), 3 ( 29 cursos ), 4 ( oito cursos ) e 5 ( um curso ).
Três cursos oferecidos por Universidades no Nordeste ficaram sem conceito ( SC ).
Os Centros Universitários contaram com dois cursos participantes na região Nordeste, um com conceito 1 e outro, com conceito 2.
As Faculdades foram representadas por 21 cursos na região Nordeste, dos quais 12 ficaram \textbf{alocados} no conceito 2, conceito modal.
Sete cursos ficaram com conceito 1 e dois com conceito 3.
Nenhum curso ficou sem conceito ( SC ) ou recebeu conceitos 4 ou 5.
Na região Sudeste, as Universidades concentraram 58 dos 120 cursos da região.
Entre os cursos em Universidades, na região, o conceito modal foi 3 com 31 cursos, e um curso ficou sem conceito ( SC ).
Os demais cursos receberam os conceitos 2 ( oito cursos ), 4 ( dez cursos ) e 5 ( oito cursos ).
Nenhum curso foi avaliado com conceito 1.
Os Centros Universitários participaram com 25 cursos na região Sudeste, dos quais 14 obtiveram conceito modal, 3.
Nenhum curso ficou sem com conceitos 1 ou 5.
Os demais receberam os conceitos 2 ( três cursos ) e 4 ( sete cursos ).
Um curso ficou sem conceito ( SC ).
As Faculdades foram representadas por 37 cursos na região Sudeste, dos quais 18 receberam conceito 3, conceito modal.
Os demais se distribuíram nos conceitos 1 ( três cursos ), 2 ( 14 cursos ), 4 ( um curso ) e 5 ( um curso ).
Nenhum curso ficou sem conceito ( SC ).
Dos 65 cursos da região Sul, 47 eram de Universidades, para os quais o conceito modal foi 3, com 20 cursos.
Nesse tipo de organização, apenas três dos cursos ficaram sem conceito ( SC ).
Os demais receberam os conceitos 1 ( dois cursos ) 2 ( cinco cursos ), 4 ( 13 cursos ) e 5 ( quatro cursos ). aos conceitos 4 ou 5 e nenhum curso ficou sem conceito ( SC ) nesta categoria e região.
Os Centros Universitários da região Norte foram representados por somente um curso, este curso foi avaliado com conceito 2.
As Faculdades participaram com cinco cursos na região Norte.
Um curso ficou sem conceito ( SC ) e outro ficou com conceito 2.
Os cursos oferecidos por Faculdades receberam conceito 3 ( três cursos ), conceito modal.
Nenhum curso foi \textbf{alocado} nos demais conceitos nesta organização e região.
Na região Nordeste, as Universidades participaram com 92 dos 115 cursos na Área de História ( Licenciatura ).
O conceito modal foi 2, com 41 cursos, enquanto os demais cursos oferecidos por Universidades receberam conceitos 1 ( dez cursos ), 3 ( 29 cursos ), 4 ( oito cursos ) e 5 ( um curso ).
Três cursos oferecidos por Universidades no Nordeste ficaram sem conceito ( SC ).
Os Centros Universitários contaram com dois cursos participantes na região Nordeste, um com conceito 1 e outro, com conceito 2.
As Faculdades foram representadas por 21 cursos na região Nordeste, dos quais 12 ficaram \textbf{alocados} no conceito 2, conceito modal.
Sete cursos ficaram com conceito 1 e dois com conceito 3.
Nenhum curso ficou sem conceito ( SC ) ou recebeu conceitos 4 ou 5.
Na região Sudeste, as Universidades concentraram 58 dos 120 cursos da região.
Entre os cursos em Universidades, na região, o conceito modal foi 3 com 31 cursos, e um curso ficou sem conceito ( SC ).
Os demais cursos receberam os conceitos 2 ( oito cursos ), 4 ( dez cursos ) e 5 ( oito cursos ).
Nenhum curso foi avaliado com conceito 1.
Os Centros Universitários participaram com 25 cursos na região Sudeste, dos quais 14 obtiveram conceito modal, 3.
Nenhum curso ficou sem com conceitos 1 ou 5.
Os demais receberam os conceitos 2 ( três cursos ) e 4 ( sete cursos ).
Um curso ficou sem conceito ( SC ).
As Faculdades foram representadas por 37 cursos na região Sudeste, dos quais 18 receberam conceito 3, conceito modal.
Os demais se distribuíram nos conceitos 1 ( três cursos ), 2 ( 14 cursos ), 4 ( um curso ) e 5 ( um curso ).
Nenhum curso ficou sem conceito ( SC ).
Dos 65 cursos da região Sul, 47 eram de Universidades, para os quais o conceito modal foi 3, com 20 cursos.
Nesse tipo de organização, apenas três dos cursos ficaram sem conceito ( SC ).
Os demais receberam os conceitos 1 ( dois cursos ) 2 ( cinco cursos ), 4 ( 13 cursos ) e 5 ( quatro cursos ).
Os Centros Universitários da região Sul ofereceram sete dos 65 cursos participantes na região.
Destes, três cursos foram \textbf{alocados} no conceito 3 ( conceito modal ), enquanto apenas um curso recebeu conceito 2, e dois cursos obtiveram conceito 4.
Nenhum curso obteve conceitos 1 ou 5.
Somente um curso ficou sem conceito ( SC ).
Foram 11 os cursos vinculados a Faculdades na região Sul, e cinco receberam conceito 3, o conceito modal.
Três cursos receberam conceito 2 e outros três cursos receberam o conceito 4.
Nenhum curso ficou sem conceito ou recebeu os conceitos 1 ou 5 nessa categoria e região.
Na região Centro-Oeste, 33 dos 42 cursos eram de Universidades.
Nesse tipo de organização, um curso ficou sem conceito ( SC ) e nenhum curso ficou com conceito 5.
O conceito modal foi 2, com 16 cursos.
Os demais receberam os conceitos 1 ( dois cursos ), 3 ( nove cursos ) e 4 ( cinco cursos ).
Os Centros Universitários da região Centro-Oeste contaram com apenas um curso, e este curso foi avaliado com conceito 3.
Dos oito cursos oferecidos em Faculdades na região Centro-Oeste, quatro ficaram \textbf{alocados} no conceito 2, conceito modal, enquanto três cursos receberam conceito 3.
Um curso ficou sem conceito ( SC ) e nenhum curso foi avaliado nos demais conceitos restantes nesta organização e região.
O Gráfico 5.6 apresenta a distribuição do Conceito Enade segundo a Organização Acadêmica da IES.
Os cursos em Centros Universitários ( linha verde ) apresentam uma poligonal mais à direita do que as demais, denotando uma melhor distribuição de Conceitos Enade.
No outro extremo, temos os cursos oferecidos em Faculdades ( linha vermelha ) com a poligonal mais à esquerda e, consequentemente, uma distribuição de Conceitos Enade com menores valores.
IES de Universidades ( linha azul ) apresentam uma situação intermediária.

% matched lemmas: alocar
\end{textsample}
